\section{Common estimates and variational tools}\label{sec:tools}
We use the continuous-dual framework and positivity criterion of
\cite[Lemma 3.3 and Appendix A]{LiZou}. The trace and variation estimates
below develop \cite[Lemma 3.12 and Proposition 3.13]{LiZou} for varying PSD
weights and fixed matrix directions.

\begin{lemma}[Continuous representation and positivity]\label{lem:duality}
For every symmetric-mass path,
\begin{equation}\label{eq:continuousdual}
 \J_{L,H}(\mu)=
 \sup_{(A,B,C)\in C(Q,\Oset_{L,H})}
 \left(\int A\cdot dG+\int B\cdot dq+\int C\cdot dR\right).
\end{equation}
Finite action forces $g\succeq0$ almost everywhere, hence $G(S)\succeq0$
for every measurable $S$. The fixed-weight action is weak-star lower
semicontinuous.
\end{lemma}
\begin{proof}
The set $\Oset_{L,H}$ contains zero and is closed because its defining
matrix expression is continuous and the cone of negative semidefinite
matrices is closed. To verify convexity, for $0\le\theta\le1$ write
$B_\theta=\theta B_1+(1-\theta)B_2$. For every $K\succeq0$,
\[
 \begin{aligned}
 &\theta B_1KB_1^T+(1-\theta)B_2KB_2^T-B_\theta KB_\theta^T\\
 &\qquad=\theta(1-\theta)(B_1-B_2)K(B_1-B_2)^T\succeq0.
 \end{aligned}
\]
Use $K=LL^T$ and the corresponding identity for $C$ with $K=HH^T$.
The constraint at a convex combination of two dual triples is therefore
bounded above by the convex combination of their constraint matrices,
which is negative semidefinite.

For every continuous admissible dual field, the definition of the support
function gives $A\cdot g+B\cdot z+C\cdot r\le j_{L,H}(g,z,r)$ pointwise.
The signed Radon--Nikodym pairing and integration give the upper bound in
\eqref{eq:continuousdual}.

For the reverse inequality, following \cite[Appendix A]{LiZou}, write
$v=(g,z,r)$ and embed the dual set in the full-coordinate Hilbert space
\[
 \begin{gathered}
 \mathcal E=\R^{n\times n}\times\R^{n\times k}\times\R^{n\times n},\\
 \norm{(A,B,C)}_2^2=\norm A_F^2+\norm B_F^2+\norm C_F^2.
 \end{gathered}
\]
The mass field of every admissible dual point remains symmetric.
The inner product is the sum of the three Frobenius pairings. Thus both
off-diagonal entries of each symmetric matrix are counted in its norm
and pairing, exactly as in the action. The preceding calculation establishes that the dual set is nonempty,
closed and convex.

\paragraph{Finite measurable selectors.}
Every subset of a finite-dimensional Euclidean space is second countable
and therefore separable.
Choose a countable dense sequence $(u_\ell)_{\ell\ge0}$ in
$\Oset_{L,H}$, with $u_0=0$. Continuity of a linear pairing in its dual
argument gives, for every $v$,
\[
 j_{L,H}(v)=\sup_{\ell\ge0}u_\ell\cdot v.
\]
Indeed, every fixed point of the dual set is a limit of sequence points,
so its pairing is bounded by this countable supremum. This reasoning
also applies when the supremum is infinite.

Choose finite Borel representatives of the Radon--Nikodym densities,
altering a null set if necessary, and put
$a_\ell=u_\ell\cdot v$ and $m_N=\max_{0\le\ell\le N}a_\ell$. These are
measurable scalar functions. Define $s_0=u_0$ recursively: keep $s_N$
when $a_{N+1}\le s_N\cdot v$, and otherwise take $s_{N+1}=u_{N+1}$.
The comparisons are measurable, so each $s_N$ is a measurable selector
whose values belong to the finite set $\{u_0,\ldots,u_N\}$; in particular
it is bounded. This rule retains the earlier choice in a tie.
It takes values in the dual set and satisfies $s_N\cdot v=m_N$.
Because $u_0=0$, the functions $m_N$ are nonnegative and increase
pointwise to $j_{L,H}(v)$.

\paragraph{Weighted continuous approximation.}
Let $v_a$ range over all matrix entries of the three densities, set
$w=\sum_a|v_a|$, and define $d\nu=w\,d\rho_\mu$. The coordinate densities
are integrable, so $\nu$ is a finite Borel measure and hence a Radon
measure on the compact metric space $Q$. Moreover $\norm v_2\le w$.
We give the approximation argument in this particular weighted measure.
For a Borel set $E$ and $\eta>0$, regularity supplies a compact $K\subset E$
and a relatively open $U\supset E$ with $\nu(U\setminus K)<\eta$.
The continuous separation of two disjoint closed sets supplies
$f:Q\to[0,1]$ equal to $1$ on $K$ and $0$ on $Q\setminus U$.
It follows that $\int_Q|f-\mathbf1_E|\,d\nu<\eta$.
Express $s_N$ as a finite sum $\sum_\ell b_\ell\mathbf1_{E_\ell}$ over
its measurable level sets. Applying the preceding approximation to these
sets and summing with the vector coefficients gives continuous fields $h:Q\to\mathcal E$
with arbitrarily small $\int_Q\norm{h-s_N}_2\,d\nu$.
For example, the triangle inequality bounds this integral by
$\sum_\ell\norm{b_\ell}_2
\int_Q|f_\ell-\mathbf1_{E_\ell}|\,d\nu$, and each scalar error can
be chosen separately. This also proves the usual density of continuous
fields in weighted $L^1$ after simple-function approximation.

\paragraph{Projection preserves the approximation.}
Let $P$ be nearest-point projection onto $\Oset_{L,H}$ in $\mathcal E$.
To obtain the projection, choose points whose distances from $x$ converge
to the infimum. Convexity puts their midpoints in the same set, and the
parallelogram identity bounds the squared distance between two chosen
points by twice the sum of their squared distances from $x$ minus four
times the squared infimum. This tends to zero, so the minimizing sequence
is Cauchy. Completeness and closedness give a minimizer $p$.
The same midpoint identity proves uniqueness: two distinct minimizers
would have a midpoint with strictly smaller squared distance. Comparing squared distances
along the segment from $p$ to any $z\in\Oset_{L,H}$ gives
\[
 \langle x-p,z-p\rangle\le0.
\]
For $p=P(x)$ and $q=P(y)$, apply this inequality first with $z=q$ and
then with the roles reversed and $z=p$. Adding yields
\[
 \norm{p-q}_2^2\le\langle p-q,x-y\rangle
                   \le\norm{p-q}_2\norm{x-y}_2.
\]
Thus $P$ is $1$-Lipschitz and continuous, including when $p=q$, and it
fixes every point of the dual set.

Given $\varepsilon>0$, choose the continuous approximation $h$ with
$\int_Q\norm{h-s_N}_2\,d\nu\le\varepsilon$ and put
$u_{N,\varepsilon}=P\circ h$. This is a continuous admissible field.
Since $P(s_N)=s_N$ pointwise, nonexpansiveness and $\norm v_2\le w$ give
\[
 \begin{aligned}
 \int_Q|(s_N-u_{N,\varepsilon})\cdot v|\,d\rho_\mu
 &\le\int_Q\norm{s_N-u_{N,\varepsilon}}_2\,d\nu\\
 &\le\int_Q\norm{s_N-h}_2\,d\nu\le\varepsilon.
 \end{aligned}
\]
Every integral here is well-defined: the selector is bounded, continuous
fields are bounded on compact $Q$, and all coordinate densities are
integrable. If $S$ denotes the continuous dual supremum, this estimate
implies $\int_Qm_N\,d\rho_\mu\le S+\varepsilon$ for each $N$.
Let $\varepsilon\downarrow0$ and then use monotone convergence to obtain
\[
 \J_{L,H}(\mu)=\lim_{N\to\infty}\int_Qm_N\,d\rho_\mu\le S.
\]
If the action is infinite, the same increasing selector integrals force
$S=+\infty$. This proves the reverse inequality.

For positivity, fix a point at which the symmetric density $g$ has finite
local cost. For each $v\in\R^n$ and $s\ge0$, the triple $(-svv^T,0,0)$ is
admissible and has pairing $-s v^Tgv$. Thus $v^Tgv<0$ implies
$j_{L,H}(g,z,r)=+\infty$, a contradiction. All quadratic forms are
therefore nonnegative at every such point. Finite action implies finite
local cost almost everywhere, hence $g\succeq0$ almost everywhere.
Consequently, for every measurable $S$ and every
real $v$, $v^TG(S)v=\int_S v^Tgv\,d\rho_\mu\ge0$.
Finally each continuous dual pairing is weak-star continuous in all
coordinates. A supremum of these continuous functions is lower
semicontinuous, proving the last assertion.
\end{proof}

\begin{lemma}[Trace and Frobenius variation estimates]\label{lem:rates}
Let $\mu\in\CE$ have finite action $E=\J_{L,H}(\mu)$, set
$T=\int_Q\tr(dG)$ and $m_0=\tr G_0(\Omega)$. Then
\begin{equation}\label{eq:trace}
 T=m_0+\int_Q(1-t)\tr(dR),\qquad
 0\le T\le2m_0+2\norm H_F^2 E.
\end{equation}
For any compatible rectangular $A$,
\begin{equation}\label{eq:rates}
 \norm{qA}_{\TV,F}^2\le2ET\norm{L^TA}_F^2,
 \qquad\norm{RA}_{\TV,F}^2\le2ET\norm{H^TA}_F^2.
\end{equation}
All variations in \eqref{eq:rates} are finite.
\end{lemma}
\begin{proof}
Test \eqref{eq:ce} with $(1-t)I_n$. Since $D^*I_n=0$, its left side is
$-T+\int_Q(1-t)\tr(dR)$ and its right side is $-m_0$, giving the trace
identity. Positivity from Lemma~\ref{lem:duality} gives $T\ge0$.
For real $s$, use the admissible reaction test
$C=s(1-t)I_n$, $B=0$, $A=-CHH^TC^T/2$. Pairing and integrating, then using
$(1-t)^2\le1$ and $HH^T\preceq\norm H_F^2I_n$, gives
\[
 s(T-m_0)-\tfrac12s^2\norm H_F^2T\le E.
\]
If $H\ne0$, substitution of $s=\norm H_F^{-2}$ gives
$(T/2-m_0)/\norm H_F^2\le E$, which is \eqref{eq:trace}.
If $H=0$, the inequality for both signs and arbitrarily large values of
$s$ forces $T=m_0$.

For the variation bound, work where the local cost is finite. For a
coordinate vector $e_i$ and $v\in\R^k$, use the triple with flux
$B=se_iv^T$, reaction $C=0$ and mass field
$-\tfrac{s^2}2\norm{L^Tv}^2e_ie_i^T$. Its pairing is
$s(zv)_i-s^2g_{ii}\norm{L^Tv}^2/2$. For real $x$ and $a>0$,
\[
 \begin{aligned}
 sx-\tfrac a2s^2&=\frac{x^2}{2a}
                   -\frac a2\left(s-\frac xa\right)^2,\\
 \sup_{s\in\R}\left(sx-\tfrac a2s^2\right)&=\frac{x^2}{2a}.
 \end{aligned}
\]
If $a=0$, finiteness of this supremum forces $x=0$. Apply these two cases
with $a=g_{ii}\norm{L^Tv}^2\ge0$ and $x=(zv)_i$. In both cases,
\[
 |(zv)_i|^2\le2j_{L,H}(g,z,r)\,g_{ii}\norm{L^Tv}^2.
\]
Applying this inequality to each column $a_b$ of the fixed matrix $A$
and summing over $i,b$ gives
\[
 \begin{aligned}
 \norm{zA}_F^2
 &\le2j_{L,H}(g,z,r)\Bigl(\sum_i g_{ii}\Bigr)
                            \Bigl(\sum_b\norm{L^Ta_b}^2\Bigr)\\
 &=2j_{L,H}(g,z,r)\tr(g)\norm{L^TA}_F^2.
 \end{aligned}
\]
The same rank-one test in the $C$ channel gives
$\norm{rA}_F^2\le2j_{L,H}(g,z,r)\tr(g)\norm{H^TA}_F^2$ for every
compatible $A$. In particular it controls the full reaction density.
Integrating square roots and applying Cauchy--Schwarz to the two
nonnegative functions $\sqrt{j_{L,H}}$ and $\sqrt{\tr g}$ yields
\[
 \begin{aligned}
 \int_Q\norm{zA}_F\,d\rho_\mu
 &\le\sqrt2\norm{L^TA}_F
       \left(\int_Qj_{L,H}\,d\rho_\mu\right)^{1/2}
       \left(\int_Q\tr g\,d\rho_\mu\right)^{1/2}\\
 &=\sqrt{2ET}\norm{L^TA}_F,
 \end{aligned}
\]
and the corresponding reaction inequality. The product densities are
integrable; Radon--Nikodym linearity identifies them with the densities
of $qA$ and $RA$. Formula~\eqref{eq:variation} and squaring give
\eqref{eq:rates}, including $E=0$, $T=0$, and null-weight directions.
\end{proof}

For a positive-definite fixed reaction weight, existence of a minimizer is
established in \cite[Theorem 4.2]{LiZou}. The following argument treats
bounded varying weights and singular fixed weights at finite value.

\begin{lemma}[Bounded-weight compactness]\label{lem:compactness}
If $(L_j,H_j)$ is bounded and admissible $\mu_j$ have uniformly bounded
actions, a common subsequence converges weak-star in all signed coordinates
to an admissible path. For fixed weights, every finite minimum is attained
with PSD measure-valued mass.
\end{lemma}
\begin{proof}
The uniform bounds on the weights and actions, together with
\eqref{eq:trace}, bound $T_j$. Taking coordinate vectors for $A$ in
\eqref{eq:rates} bounds the Frobenius variation of every column of $q_j$
and $R_j$, hence the variation of every scalar coordinate. Positivity
also controls all mass coordinates: the $2\times2$ principal minors of
$g_j\succeq0$ give
\[
 |(g_j)_{ab}|\le\sqrt{(g_j)_{aa}(g_j)_{bb}}
 \le\tfrac12\bigl((g_j)_{aa}+(g_j)_{bb}\bigr)\le\tr g_j.
\]
Integrating shows $|G_j^{ab}|(Q)\le T_j$. Thus every coordinate has a
uniform total-variation bound.

We spell out why these bounds give a common subsequence. If $Q$ is
empty, every coordinate measure is zero; otherwise choose a point of $Q$.
Decompose each signed coordinate $\eta_j$ into its positive and negative Jordan parts.
Both are finite positive measures with mass bounded by the common
variation bound $C$. A positive measure $\gamma$ of mass $m\le C$ can be
written as $m\pi$, where $\pi$ is a probability measure; if $m=0$, any
fixed probability measure gives the same zero measure. The space of
probability measures on compact metric $Q$ is compact and metrizable for
weak convergence. Thus $[0,C]$ times that space is sequentially compact.
There are only finitely many coordinates and two Jordan parts for each,
so compactness of their finite product supplies a single subsequence
along which all masses and all normalized positive measures converge.
Multiplication by the masses gives weak limits of all the Jordan parts.

Subtract each limiting negative part from its limiting positive part.
The resulting finite signed Radon measure $\eta$ satisfies
$\int f\,d\eta_j\to\int f\,d\eta$ for every $f\in C(Q)$, because
this pairing is the difference of the two positive-measure pairings.
The limiting positive and negative parts need not remain mutually
singular; their difference still has exactly this pairing. Applying the
construction simultaneously to all coordinates proves the asserted
weak-star convergence. The relations
$G_j^{ab}=G_j^{ba}$ pass to the limit, so the limiting mass is symmetric.
For each fixed closed-domain $C^1$ test in \eqref{eq:ce}, all coordinates
of $\partial_t\Phi$, $D^*\Phi$ and $\Phi$ are continuous on $Q$.
Every integral in that equation therefore converges, while its prescribed
endpoint terms are unchanged. The limiting path is admissible.

For fixed weights and finite infimum $V(L,H)$, choose admissible paths
with actions at most $V(L,H)+1/(j+1)$. This is a bounded finite-action
minimizing sequence, so the preceding argument gives an admissible
subsequential limit $\mu$. Lemma~\ref{lem:duality} gives
$\J_{L,H}(\mu)\le V(L,H)$, and the definition of the infimum gives the
opposite inequality. Thus the minimum is attained, and its finite action
implies PSD measure-valued mass.
\end{proof}

\begin{lemma}[Joint lower bound]\label{lem:jointlsc}
If $L_j\to L_*$, $H_j\to H_*$ and $\mu_j\weak\mu$, then
\begin{equation}\label{eq:jointlsc}
 \J_{L_*,H_*}(\mu)\le\liminf_j\J_{L_j,H_j}(\mu_j).
\end{equation}
No relative comparison is needed here.
\end{lemma}
\begin{proof}
If the right side of \eqref{eq:jointlsc} is $+\infty$, the inequality is
automatic. Otherwise pass to a subsequence whose finite actions converge to that
finite lower limit, denoted by $\ell$, and relabel it. Fix a continuous
target dual field $(A,B,C)$. We compensate for the change in its quadratic
constraint by subtracting a scalar multiple of $I_n$ from its mass field.
Let $M_B=\sup_Q\norm B_F^2$, $M_C=\sup_Q\norm C_F^2$, write
$|S|_1=\sum_{a,b}|S_{ab}|$, and set
\[
 \delta_j=\tfrac12\bigl(M_B|L_jL_j^T-L_*L_*^T|_1
                  +M_C|H_jH_j^T-H_*H_*^T|_1\bigr).
\]
For any symmetric $S$, compatible matrix $F$ and vector $x$, putting
$y=F^Tx$ gives
\[
 |x^TFSF^Tx|
 \le\sum_{a,b}|S_{ab}|\,|y_a|\,|y_b|
 \le |S|_1\norm F_F^2|x|^2.
\]
Consequently $FSF^T\preceq |S|_1\norm F_F^2I_n$. Apply this to the two
Gram differences. The change in the quadratic constraint is bounded
above by $\delta_jI_n$, uniformly on $Q$, so
$(A-\delta_jI_n,B,C)$ is admissible at $(L_j,H_j)$. Matrix convergence
implies $\delta_j\to0$.

The corrected field has pairing
\[
 \int_Q A\cdot dG_j+\int_Q B\cdot dq_j+\int_Q C\cdot dR_j
       -\delta_j\int_Q\tr(dG_j).
\]
Weak-star convergence makes the first three terms converge to the target
pairing. It also makes the trace pairing converge, since it is a finite
sum of coordinate pairings against the constant function $1$; thus the
last term tends to zero. This step uses no positivity assumption on the
limiting mass. The defining support inequality bounds the corrected
pairing by $\J_{L_j,H_j}(\mu_j)$, so its limit is at most $\ell$.
Every continuous target dual pairing is therefore at most $\ell$.
Taking their supremum in \eqref{eq:continuousdual} proves
$\J_{L_*,H_*}(\mu)\le\ell$, including the assertion that a finite
lower limit cannot coexist with infinite target action.
\end{proof}

\begin{lemma}[Relative Gram comparison]\label{lem:comparison}
If $\alpha>0$ and
\begin{equation}\label{eq:relative}
 LL^T\succeq\alpha L_*L_*^T,
 \qquad HH^T\succeq\alpha H_*H_*^T,
\end{equation}
then for every path $\mu$,
\begin{equation}\label{eq:comparison}
 \J_{L,H}(\mu)\le\alpha^{-1}\J_{L_*,H_*}(\mu).
\end{equation}
\end{lemma}
\begin{proof}
For $(A,B,C)\in\Oset_{L,H}$, the target slack of
$(\alpha A,\alpha B,\alpha C)$ is
\[
 \alpha\bigl(-A-\tfrac12BLL^TB^T-\tfrac12CHH^TC^T\bigr)
 +\tfrac\alpha2 B(LL^T-\alpha L_*L_*^T)B^T
 +\tfrac\alpha2 C(HH^T-\alpha H_*H_*^T)C^T.
\]
Each term is PSD, so multiplication by $\alpha$ sends the current dual
set into the target dual set and multiplies its pairing by $\alpha$.
Taking support suprema and integrating proves \eqref{eq:comparison},
including infinite actions.
\end{proof}
